\documentclass[a4,11pt]{aleph-notas}

% -- Paquetes adicionales
\usepackage{enumitem}
\usepackage{aleph-comandos}

% -- Datos
\institucion{Facultad de Ciencias Exactas, Naturales y Ambientales}
\carrera{Catálogo STEM}
\asignatura{Matemática III}
\tema{Resumen 02: Geometría de curvas}
\autor[A. Merino]{Andrés Merino}
\fecha{Periodo 2026-1}

\logouno[0.16\textwidth]{Logos/logoPUCE_04_ac}
\definecolor{colortext}{HTML}{1957d1}
\definecolor{colordef}{HTML}{1957d1}
\fuente{montserrat}

\begin{document}

\encabezado

%%%%%%%%%%%%%%%%%%%%%%%%%%%%%%%%%%%%%%
\section{Vectores y planos asociados}
%%%%%%%%%%%%%%%%%%%%%%%%%%%%%%%%%%%%%%

En esta sección consideraremos $\func{\alpha}{I}{\R^3}$, dos veces derivable tal que $\alpha'(t)\neq 0$ para todo $t\in I$.

\begin{defi}[Vector unitario tangente]
    Se define el vector unitario tangente por
    \[
        T_\alpha(t)=\frac{\alpha'(t)}{\|\alpha'(t)\|},
    \]
    para todo $t\in I$.
\end{defi}

\begin{defi}[Vector normal principal]
    Si $T_\alpha'(t)\neq 0$ para todo $t\in I$, se define el vector normal principal por
    \[
        N_\alpha(t)=\frac{T_\alpha'(t)}{\|T_\alpha'(t)\|},
    \]
    para todo $t\in I$.
\end{defi}

\begin{defi}[Vector binormal principal]
    Si $T_\alpha'(t)\neq 0$ para todo $t\in I$, se define el vector binormal principal por
    \[
        B_\alpha(t)=T_\alpha(t)\times N_\alpha(t),
    \]
    para todo $t\in I$.
\end{defi}

\begin{defi}[Plano osculador]
    El plano
    \[
        P(\alpha(t);T_\alpha(t),N_\alpha(t))
    \]
    se lo denomina plano osculador de la trayectoria en el punto $\alpha(t)$.
\end{defi}

\begin{defi}[Plano normal]
    El plano
    \[
        P(\alpha(t);B_\alpha(t),N_\alpha(t))
    \]
    se lo denomina plano normal de la trayectoria en el punto $\alpha(t)$.
\end{defi}

\begin{defi}[Plano rectificante]
    El plano
    \[
        P(\alpha(t);T_\alpha(t),B_\alpha(t))
    \]
    se lo denomina plano rectificante de la trayectoria en el punto $\alpha(t)$.
\end{defi}

\begin{prop}
    El vector aceleración en el tiempo $t$ es una combinación lineal de los vectores $T_\alpha(t)$ y $T_\alpha'(t)$ para todo $t\in I$, además
    \[
        a_\alpha(t)=\rho'_\alpha(t) T_\alpha(t)+\rho_\alpha(t) T_\alpha'(t).
    \]
\end{prop}

\begin{advertencia}
    En las definiciones anteriores, si no existe ambigüedad en cuanto a la parametrización de la curva $C$, se omitirá el uso del subíndice $\alpha$.
\end{advertencia}

%%%%%%%%%%%%%%%%%%%%%%%%%%%%%%%%%%%%%%
\section{Longitud de arco y curvatura}
%%%%%%%%%%%%%%%%%%%%%%%%%%%%%%%%%%%%%%

\begin{defi}[Longitud de arco]
    Si $\func{\alpha}{[a,b]}{\R^n}$, la longitud de arco de la curva $C$ está dada por
    \[
        \ell(C)=\int_a^b \rho_\alpha(t)\, dt.
    \]
\end{defi}

\begin{prop}
    La longitud de arco es invariante bajo cambio de parámetro.
\end{prop}

En esta sección consideraremos $\func{\alpha}{I}{\R^3}$, dos veces derivable, con $\alpha'(t)\neq 0$ para todo $t\in I$.

\begin{defi}[Curvatura]
    El vector curvatura se define por
    \[
        K_\alpha(t)=\frac{1}{\rho_\alpha(t)}T_\alpha'(t),
    \]
    para todo $t\in I$. Además, la curvatura se define por
    \[
        \kappa_\alpha(t)=\frac{\|T_\alpha'(t)\|}{\rho_\alpha(t)},
    \]
    para todo $t\in I$.
\end{defi}

\begin{prop}
    La curvatura es invariante bajo cambio de parámetro.
\end{prop}

\begin{prop}
    Se tiene que
    \[
        a_\alpha(t)=\rho_\alpha'(t)\, T_\alpha(t) + \kappa_\alpha(t)\, \rho^2_\alpha(t)\, N_\alpha(t),
    \]
    para todo $t\in I$. Además,
    \[
        \kappa_\alpha(t) = \frac{\|a_\alpha(t)\times v_\alpha(t)\|}{\rho^3_\alpha(t)}
    \]
    para todo $t\in I$.
\end{prop}

\begin{defi}[Torsión]
    La torsión se define por
    \[
        \tau_\alpha(t)= \frac{v_\alpha(t)\cdot(a_\alpha(t)\times a_\alpha'(t))}{\| v_\alpha(t)\times a_\alpha(t) \|^2},
    \]
    para todo $t\in I$.
\end{defi}

\begin{prop}
    La torsión es invariante bajo cambio de parámetro.
\end{prop}

\end{document}
